\chapter{Inserting a Pacemaker}

\section*{Definition and Overview}
A pacemaker is a small electronic device implanted under the skin of the chest that monitors the heart and delivers tiny electrical impulses to keep the heart rate from falling too low. It is used to treat bradycardia, an abnormally slow heart rhythm that causes symptoms such as dizziness, fainting, breathlessness, and fatigue. The device consists of a pulse generator and one or more leads (wires) passed through a vein into the heart. Insertion is performed under local anaesthetic with sedation, usually as a day or overnight procedure, and takes about an hour. Pacemakers are safe, reliable, and life-changing for most recipients, and modern devices are small, long-lasting, and compatible with normal life.

\section*{Epidemiology}
Pacemakers are implanted in many thousands of Australians each year, and the number increases with age, with most recipients over 70. The most common indication is symptomatic bradycardia due to sick sinus syndrome or heart block, conditions in which the heart's natural electrical system fails. Pacemakers are also used in some people with heart failure to coordinate the heartbeat (cardiac resynchronisation therapy). As the population ages, the rate of implantation continues to grow. Pacemaker insertion is a routine procedure with low complication rates in experienced centres.

\section*{Aetiology and Pathophysiology}
The heart's rhythm is set by the sinus node, a natural pacemaker in the right upper chamber, and conducted through the atrioventricular node to the ventricles. Disease or degeneration of these structures, from ageing, ischaemic heart disease, heart attacks, inflammation, medicines, or heart surgery, can slow or block conduction. When the heart rate falls below what the body needs, symptoms develop, and some rhythms can be life-threatening. A pacemaker replaces or supports the failing system: the generator senses the heart's own activity and, when a beat is missed or too slow, delivers a stimulus that triggers a contraction. In rate-responsive models, sensors increase the rate with activity.

\section*{Clinical Features}
\begin{itemize}
    \item Fainting or near-fainting (syncope) from slow heart rhythm
    \item Dizziness, light-headedness, and confusion
    \item Breathlessness on exertion and reduced exercise tolerance
    \item Fatigue and tiredness
    \item Palpitations, or a sense that the heart is skipping or pausing
    \item Symptoms may be intermittent, making diagnosis challenging
    \item After implantation: a small wound over the collarbone, with restriction of arm movement for a few weeks
    \item The device is usually invisible under the skin and can be checked remotely
\end{itemize}

\section*{Differential Diagnosis}
Symptoms that suggest slow heart rhythm may have other causes.
\begin{itemize}
    \item \textbf{Vasovagal syncope:} fainting from a drop in blood pressure, not slow rhythm
    \item \textbf{Orthostatic hypotension:} blood pressure fall on standing
    \item \textbf{Atrial fibrillation and other arrhythmias:} fast or irregular rhythms
    \item \textbf{Medicines:} beta-blockers, digoxin, and others slowing the heart
    \item \textbf{Thyroid disease and anaemia:} causing fatigue and breathlessness
    \item \textbf{Seizures:} may mimic fainting
    \item \textbf{Structural heart disease:} valve or muscle disease
\end{itemize}

\section*{When to Seek Medical Care}
Anyone with unexplained fainting, dizziness, or breathlessness should be assessed, and investigation with ECG monitoring may be needed. People with a pacemaker should attend scheduled checks and seek review for wound problems, unusual symptoms, or an alert from their device. Pacemaker checks are usually done in person or remotely every few months.

\textbf{Call 000 (or local emergency services) for:}
\begin{itemize}
    \item Fainting, collapse, or severe dizziness
    \item Chest pain, breathlessness, or palpitations
    \item Signs of infection at the pacemaker site: redness, swelling, pain, or fever
    \item Wound dehiscence or bleeding that will not stop
    \item A sudden shock or repeated shocks from the device
    \item Symptoms of stroke: face, arm, or leg weakness, or difficulty speaking
\end{itemize}

\section*{Diagnosis and Investigations}
\begin{itemize}
    \item \textbf{Electrocardiogram (ECG):} detects slow rhythms and heart block
    \item \textbf{Holter or event monitoring:} prolonged recording to catch intermittent bradycardia
    \item \textbf{Echocardiogram:} assessment of heart structure and function
    \item \textbf{Exercise testing:} where symptoms occur with exertion
    \item \textbf{Blood tests:} thyroid function and electrolyte levels
    \item \textbf{Assessment of medicines:} identifying drugs that slow the heart
    \item \textbf{Pre-procedure review:} blood tests, medicines, and fitness for the procedure
\end{itemize}

\section*{Management}
\begin{itemize}
    \item \textbf{Insertion procedure:} the generator is placed under the skin and leads passed through a vein into the heart, under X-ray guidance
    \item \textbf{Local anaesthesia with sedation:} the person is awake but comfortable
    \item \textbf{Device programming:} the pacemaker is set for the individual's needs and checked before discharge
    \item \textbf{Aftercare:} wound care, arm restrictions for several weeks, and simple analgesia
    \item \textbf{Follow-up:} device checks in clinic or remotely, with battery monitoring
    \item \textbf{Medication review:} adjusting medicines that affect the heart rate
    \item \textbf{Addressing underlying heart disease:} treatment of ischaemia, heart failure, and valve disease
    \item \textbf{Lifestyle guidance:} most activities are safe, with care around heavy lifting and strong magnetic fields
\end{itemize}

\section*{Complications}
\begin{itemize}
    \item Infection of the wound or device, requiring antibiotics and sometimes removal
    \item Bleeding, bruising, or a pocket of blood at the site
    \item Lead problems: displacement, fracture, or damage to the heart or blood vessels
    \item Pneumothorax, a collapsed lung from the vein puncture
    \item Rarely, perforation of the heart wall
    \item Device-related issues: battery depletion, sensing problems, and inappropriate activity
    \item Restrictions on some activities, including contact sport and strong electromagnetic fields
\end{itemize}

\section*{Prognosis and Outcomes}
Pacemakers dramatically improve the quality of life and safety of people with symptomatic bradycardia, and most recipients are free of dizziness and fainting after implantation. Modern devices last 8 to 12 years before the battery needs replacing, and follow-up is simple, often remote. The main risks are procedural and early, and long-term reliability is excellent. People with pacemakers can expect to live normal, active lives, including work, travel, and most sports.

\section*{Prevention and Self-Care}
\begin{itemize}
    \item Attend all device checks, in person or remotely
    \item Keep the wound dry until healed and report redness, swelling, or pain
    \item Avoid heavy lifting and vigorous arm movement on the device side for several weeks
    \item Carry your pacemaker identification card
    \item Tell all doctors and dentists about your pacemaker
    \item Avoid strong magnetic fields, including some MRI unless the device is compatible
    \item Keep mobile phones away from the device site
    \item Follow advice about contact sport and workplace equipment
    \item Report fainting, palpitations, or shocks promptly
\end{itemize}

\section*{Sources and Further Reading}
The following reputable sources provide further information for healthcare students and patients seeking reliable, current advice about pacemaker insertion.
\begin{itemize}
    \item Heart Foundation Australia. \textit{Pacemakers and heart rhythm.} https://www.heartfoundation.org.au/
    \item Healthdirect Australia. \textit{Pacemaker insertion.} https://www.healthdirect.gov.au/
    \item Cardiac Society of Australia and New Zealand. \textit{Device therapy information.} https://www.csanz.edu.au/
    \item Royal Australasian College of Surgeons. \textit{Pacemaker procedure information.} https://www.surgeons.org/
    \item NHS. \textit{Pacemaker implantation.} https://www.nhs.uk/conditions/pacemaker/
    \item Mayo Clinic. \textit{Pacemaker.} https://www.mayoclinic.org/
\end{itemize}
